\section{Related Work}

\paragraph{Retrieval-augmented generation.}
Retrieval-augmented generation connects parametric language models to external
evidence at inference time \citep{lewis2020rag,karpukhin2020dpr}. The standard
pipeline retrieves passages, optionally reranks or compresses them, and then
conditions a generator on the selected context
\citep{nogueira2019passage,glass2022re2g,xu2024recomp,izacard2022atlas}.
Recent surveys document many variants of this pattern, including adaptive
retrieval, query rewriting, reranking, and retrieval-time filtering
\citep{gao2024ragsurvey,chen2024benchmarking}. Much of this work reports
retrieval quality and answer quality side by side. That convention is useful,
but it leaves a gap: a high-scoring passage set can still omit, fragment, or
contradict the evidence needed to support the generated answer.

\paragraph{Retrieval evaluation versus answer support.}
Dense retrievers and cross-encoders are normally optimized for local
query--passage relevance \citep{karpukhin2020dpr,reimers2019sbert,he2021deberta}.
Benchmarks such as BEIR made cross-domain retrieval transfer a first-class
evaluation problem \citep{thakur2021beir}. RAG faithfulness is a different
estimand. The generator observes a set of passages, not an isolated passage,
and the answer may depend on coverage, redundancy, ordering, contradiction,
and whether the needed span is present at all. Work on distracting context and
long-context position effects shows that extra retrieved text can change model
behavior even when some relevant evidence is present
\citep{shi2023large,liu2024lostmiddle}. \method{} brings that distinction into
the evaluation protocol by separating local similarity from set-level
structure and answer support.

\paragraph{Faithfulness, factuality, and hallucination metrics.}
Automatic faithfulness evaluation has a long record of useful but imperfect
proxies. NLI-style factual consistency metrics, summarization consistency
benchmarks, RAGAS-style RAG metrics, and atomic factuality metrics all capture
different slices of the problem
\citep{maynez2020faithfulness,honovich2022true,laban2022summac,es2024ragas,min2023factscore}.
Surveys of hallucination emphasize the same point: unsupported generation is
not a single observable variable \citep{ji2023hallucination}. Our contribution
is not another faithfulness metric. We hold the generations fixed and show that
three plausible faithfulness observers yield materially different intervention
magnitudes.

\paragraph{LLM-as-judge and evaluator disagreement.}
LLM-based judges have become common because they are flexible and inexpensive
relative to full human annotation
\citep{liu2023geval,zheng2023judge}. They also introduce evaluator choice as
a modeling decision. Prior work studies human agreement, preference
benchmarks, and judge bias in open-ended generation. In RAG, the same issue
appears as an identification problem: if a retrieval intervention changes
DeBERTa faithfulness by 1.1 points but a local RAGAS-style judge by 14.0
points on the same generations, the intervention effect is not identified by a
single metric. The disagreement is part of the claim being estimated.

\paragraph{Corrective and self-reflective RAG.}
Corrective RAG methods try to detect retrieval failure and repair the evidence
before generation, while Self-RAG trains a model to retrieve, generate, and
critique using reflection tokens \citep{yan2024crag,asai2024selfrag}. These
methods are natural baselines for our audit because they explicitly respond to
retrieval failure rather than only changing the retriever. Our results should
not be read as a global ranking of these systems. They show that in a matched
short-answer harness, policy comparisons change across SQuAD and HotpotQA and
depend on hallucination rate, mean latency, p99 latency, and indexing cost.

\paragraph{Hierarchical retrieval and cost accounting.}
RAPTOR builds tree-structured summaries so retrieval can operate over both
raw chunks and higher-level summaries \citep{raptor2024}. That design can help
when answer support is distributed across passages, but it also introduces an
offline tree-building cost. Many RAG papers report online answer quality while
leaving indexing cost and tail latency outside the main result table.
\method{} treats those quantities as part of the deployment claim. A method
that improves faithfulness on one dataset but has larger p99 latency or higher
offline cost has not established deployment superiority by faithfulness alone.

\paragraph{Benchmarks and cross-dataset transfer.}
Cross-dataset evaluation is standard in retrieval, but less consistently used
for RAG thresholds and faithfulness interventions. BEIR-style transfer asks
whether a retriever trained or selected in one domain remains effective in
another \citep{thakur2021beir}. We ask an analogous question for RAG policies:
does a threshold tuned to recover faithfulness on one dataset transfer to the
others? The answer is uneven in our audit, with SQuAD, PubMedQA, and
NaturalQuestions showing large in-domain versus off-diagonal gaps.

\paragraph{What prior work leaves open.}
Prior work evaluates RAG systems, retrieval models, and faithfulness metrics.
This paper instead tests whether common RAG faithfulness claims are
identifiable once local similarity, context-set coherence, evaluator choice,
threshold transfer, and cost are separated. That shift is the main difference:
we are not proposing a new dominant RAG policy, but a controlled protocol for
deciding which claims the evidence supports.
